\subsection{Deliverable 7: Notebook Reflection Questions}
\label{sec:deliverable7}

\textbf{Important notes:}
\begin{itemize}
  \item You will write \textasciitilde20-25\% code (basic exploration, running analyses, creating visualisations)
  \item Focus is on \textbf{choosing appropriate methods}, \textbf{interpreting results}, and \textbf{critical analysis}
  \item \textbf{All numbered "Reflection Questions" should be answered in your main report (Deliverable 7), not in this notebook}
  \item There is no single "correct" set of methods --- you must justify your choices
  \item In-notebook analyses are brief; deeper reflections go in your report
\end{itemize}

\textbf{Part 3: Method Selection and Justification}

\textbf{Task :} From the methods above (Methods 1-11), choose 4-5 methods to run.

Complete this table in your main report.

\begin{table}[H]
\centering
\small
\label{tab:d7_method_selection}
\rowcolors{3}{black!4}{white}
\begin{tabularx}{\textwidth}{
>{\raggedright\arraybackslash}p{0.13\textwidth}
>{\raggedright\arraybackslash}p{0.14\textwidth}
>{\raggedright\arraybackslash}p{0.14\textwidth}
>{\raggedright\arraybackslash}p{0.13\textwidth}
>{\raggedright\arraybackslash}X
>{\raggedright\arraybackslash}p{0.15\textwidth}
}
\toprule
Method & DQ Dimension & Primary Facets & Relevant to RQ(s) & Why did you choose this? & Expected findings \\
\midrule
Example: Method 1 &  &  &  &  &  \\
\bottomrule
\end{tabularx}
\rowcolors{3}{}{}
\end{table}

Why did you NOT choose the other methods? What trade-offs did you consider? How did the three research questions (RQs) influence your method selection?

\textbf{Answer:}

\textbf{→ Reflection Question 1} (answer in your main report):

What kind of data is logged in the file? Assuming data quality is good, how can the data be used to answer the RQs?

\textbf{Answer:}

\textbf{→ Reflection Question 3} (answer in your main report):

Provide detailed interpretation of results from each chosen method:
\begin{itemize}
  \item Include key results.
  \item Describe and interpret the results. Did they work as expected? Did they reveal any surprising insights? Could they be improved?
  \item Were your expectations correct?
  \item Which facet(s) does this result primarily inform about?
\end{itemize}

\textbf{Answer:}

\textbf{→ Reflection Question 4} (answer in your main report):

Expand on your proposed DQ dimension.
\begin{itemize}
  \item Define your dimension.
  \item \textit{What specific pattern or issue in the data motivated this dimension? Reference your exploration code above. Provide concrete examples.}
  \item \textit{Which research question(s) would be affected by this quality issue? How?}
  \item Map the dimension relevancy in across the five facets
\end{itemize}

\textbf{Facet Mapping:}

\begin{table}[H]
\centering
\small
\label{tab:d7_facet_mapping}
\rowcolors{3}{black!4}{white}
\begin{tabularx}{\textwidth}{l c X}
\toprule
Facet & Involvement (++, +, -) & Justification \\
\midrule
Data &  &  \\
Source &  &  \\
System &  &  \\
Task &  &  \\
Human &  &  \\
\bottomrule
\end{tabularx}
\rowcolors{3}{}{}
\end{table}

\textbf{Answer:}

\textbf{→ Reflection Question 5} (answer in your main report):

For each of the five data quality facets, interpret the results you got when running selected methods and assess the quality of the dataset. You can use these questions as guidance,
\begin{itemize}
  \item What did your methods reveal about data quality at [a facet] level?
  \item What are the main data quality issues?
  \item What could be the consequences in regard to the RQs?
\end{itemize}

\textbf{Then discuss:}

a) Which facet was most challenging to assess with the available methods? Why?\\
b) How else could these data quality facets be assessed?

\textbf{Answer:}

\textbf{7.1 Invisible Work Analysis}

Complete the table with at least \textbf{3 activities} from the workflow diagram that do NOT appear in the log:

\begin{table}[H]
\centering
\small
\label{tab:d7_invisible_work}
\rowcolors{3}{black!4}{white}
\begin{tabularx}{\textwidth}{
>{\raggedright\arraybackslash}p{0.23\textwidth}
>{\raggedright\arraybackslash}p{0.17\textwidth}
>{\raggedright\arraybackslash}p{0.18\textwidth}
>{\raggedright\arraybackslash}p{0.20\textwidth}
>{\raggedright\arraybackslash}X
}
\toprule
Activity from Diagram & Why not logged? & Which facet explains this? & Consequences & How could it be captured? \\
\midrule
 &  &  &  &  \\
 &  &  &  &  \\
 &  &  &  &  \\
\bottomrule
\end{tabularx}
\rowcolors{3}{}{}
\end{table}

\textbf{→ Reflection Question 6} (answer in your main report):

a) For each activity, explain why it might not be logged (technical, social, economic reasons)\\
b) Reflect on how a process model built \emph{only} from the event log would fail to capture these activities. What are the potential consequences for system design, in regard to each of the research questions?\\
c) How does your Part I experience (interview vs. observation) parallel the gap between event log and workflow diagram?\\
d) What does the gap between logged data and actual practice tell us about building healthcare IT systems from ``data-driven'' approaches alone?\\
e) Propose one concrete change to the logging system that would capture currently invisible work. Draw on concepts from Chapter 9 (Formal and informal information systems) to discuss trade-offs: Why is complete logging difficult or even unwelcome in real-world settings? How could negative impacts be mitigated?

\textbf{Answer:}
